% Discussion section: findings F12, F17--F25 (scale gap, parity regime change, certification, PMA dose-response, auction selection, closed-loop inversion, conversion damage)

The headline result reverses a seeming architectural failure: a weight-tied single transformer block failed to match a 12-layer explicit model at 121M parameters under implicit-differentiation training (0.787 BLiMP ratio, F20), yet reaches parity (0.991 ratio) when trained with unrolled anytime supervision (F24). This is not a statement about the weight-tied architecture itself, but about the training regime. Under implicit training, the fixed-point solve becomes a bottleneck: at fixed iteration budget (12 iterations), the map's effective contraction degrades as width grows, widening the quality gap from 7\% to 21\% (F13 to F20). Unrolled training removes that bottleneck by supplying gradients directly to intermediate iterates $z_4, z_8, z_{12}$, allowing the model to learn both capacity and contraction simultaneously. The architecture's depth-memory advantage persists---serving-time $O(1)$ memory versus explicit's $O(L)$---and the anytime-trained checkpoint exhibits the pre-registered budget dial (graceful degradation across solver budgets $\{4,8,12\}$, F24).

A subsequent audit qualifies that parity in a way that changes what it recommends. The budget held equal was parameters and iteration count, not arithmetic. Matching 121M parameters with one tied block rather than twelve distinct ones forces the block to be wide, and the two models end up inversely composed: the tied model spends 85.9M parameters on embeddings and 34.8M on its block, the explicit model 38.7M and 85.1M. Width is quadratic in the dominant operation, so the tied block costs $4.92\times$ an explicit layer per token and twelve iterations cost $4.92\times$ twelve layers. Equal arithmetic is reached at $2.44$ iterations, where the ratio is $0.72$. The practitioner-facing conclusion is therefore not that equilibrium depth is free, but that weight tying trades arithmetic for parameters at a measurable exchange rate, and is worth making where parameters are the binding constraint (F44). The lesson for practitioners is practical: implicit-differentiation training of deep-equilibrium models requires tighter solvers or architectural changes; unrolled supervision with intermediate supervision is a working regime at this scale.

Quality and certification emerged as separable problems (F24). Arm B1 (anytime training) achieves 0.991 parity with the explicit baseline---the quality half of the H6 gate---at zero convergence rate (the solver still uses all 12 iterations). Arm B2 applies a finite-difference trajectory-local contraction penalty (probing only on the solve ray, not globally) and achieves the program's first certified 121M equilibrium: convergence rate 1.0 at 4.0 mean iterations, lowest memory cost (5.5 GB). Yet B2's quality drops to 0.577 BLiMP. The naive prediction, pre-registered in arm B4, was that combining anytime supervision with a raw penalty would satisfy both criteria. Experiment refuted this: the unnormalized penalty ($3.5 \times 10^4$ at initialization) drowned the cross-entropy supervision under gradient clipping, scoring 0.529, below control. The failure points to a scale mismatch: the penalty scale must track the supervision signal magnitude, a log-scale reformulation being the natural next candidate. This open problem is now precisely posed---not ``how to certify an LM,'' but ``how to preserve quality while achieving certification''---with evidence that the two objectives require distinct optimization paths at this width.

Magnetic preference optimization (PMA, parameter-space magnetic anchoring) is real but second-order to the DPO gradient (F21). Measuring through an exact preference signal (BLiMP minimal pairs, no reward model), the magnetic pull was applied inside MagneticAdamW at $\tau \in \{10^{-3}, 10^{-2}\}$, where Sokota et al.'s theory predicts it should help. KL trajectories diverged across arms (ruling out a bug hypothesis), yet held-out accuracy matched DPO \emph{exactly} across seeds and bases, and KL reduction was insignificant. A rider dose-response sweep ($\tau \in \{0.1, 1, 10\}$) refuted its own pre-registered prediction that visible drift reduction would appear by $\tau = 1$: even at $\tau = 10$, KL moved only $1.241 \to 1.226$ nats/token ($-1.2\%$) with accuracy flat. The displacement scales as $\eta\tau T \approx 2\times 10^{-5}$ relative, under-dosed for the preference signal. This is distinct from policy-space Magnetic Preference Optimization (Wang et al.\ 2025, \cite{Wang2025MPO}), which applies magnetic mirror descent in the policy distribution with self-play to target the Nash equilibrium of a preference game---a different mechanism operating on a different object. The finding for parameter-space PMA is sharp: the magnetic anchor is stable but inert at preference-tuning scales. A secondary finding compounds this: equilibrium architectures are intrinsically drift-resistant (EqLM moved KL only 0.00075 nats/token under identical loss/rate/steps, 1655$\times$ less than the explicit baseline), a double-edged property offering stability at the cost of preference insensitivity. The pretraining finding (F12)---that the magnet is loss-neutral even at $\tau = 10^{-2}$ during pretraining---emerges as consistent: the magnet's natural home is regularized Nash dynamics (periodic reference resets), not preference-phase drift control.

Token-auction selection is a scoring-time phenomenon, not an end-to-end property. At scoring time, the auction beats the best single specialist by 23\% perplexity (F22), selecting per-token via confidence bids that target independent valuations. This advantage inverted in closed-loop generation (F23): the best single specialist (NLL 3.39--3.69 tokens) defeats the auction (4.74--4.60), with the auction drifting toward one specialist's style regardless of prompt domain. Exposure bias (Bengio et al., \cite{Bengio2015}) offers the mechanistic explanation: teacher forcing supplies the true context at every decode step, re-anchoring selection to the appropriate specialist; autoregressive generation breaks the anchor, and sequential choice-compounding amplifies style drift. A separate finding, orthogonal to the verdict, is that bid competition suppresses degeneration: the auction exhibits the lowest 3-gram repetition of all systems (0.39--0.48) while uniform logit averaging collapses (0.78--0.83), consistent with the analysis of Holtzman et al.\ (\cite{Holtzman2020}) that averaging makes model collapse more likely. This anti-repetition property, decoupled from the auction's domain-selection failure, may be the mechanism worth isolating in future work. Context-aware bids (bidding on the prefix's detected domain rather than bare confidence) and tournament-style specialist selection for generation are the identified follow-ups.

The closed-loop evaluation exposed a metric pathology: the judge (the 124M BabyLM-trained explicit baseline) carries a strong style prior for child-directed speech that dwarfs between-system effects. This was discovered via a Tarka-required per-domain decomposition: the child-speech specialist scores best even on Wikipedia-only prompts (3.27--3.78 NLL), while the Wikipedia specialist generating its own register on its own domain scores 5.10--5.49---the judge's $\approx$2-nat prior overwhelms the actual domain selectivity of the systems. Accordingly, the F23 verdict is scoped as judge-relative. Two observations survive independently: the auction's output measurably diverges from teacher-forced behavior (domain consistency shifts from near-unity on in-domain prompts to near-zero on out-of-domain ones), and the repetition advantage is real and judge-free. The F22 scoping (teacher-forced scoring-time selection, not autoregressive generation) is thereby empirically vindicated---a distinction the adversarial review forced into the claim's scope before the follow-up experiment proved its necessity.

For practitioners converting pretrained stacks to recursive topologies, the conversion damage curve (F25) offers guidance. On Qwen3-1.7B, all configurations start heavily damaged (300--3000$\times$ base perplexity) before uptraining. The damage landscape is shaped by two effects: (a) average-initialization of tied layers preserves far more function than stepwise adoption (10--100$\times$ better for the same number of cores), settling an open choice in recursive-uptraining recipes; and (b) explicit outer layers matter more than the number of recursive cores (going from 1 to 2 cores at fixed outer depth changes ppl 1909 to 1933, a $\sim$1\% difference, while adding two explicit outer layers improves ppl 10-fold). The selected operating point (8 pre + 8 post explicit layers with 1 recursive core, average init, 1.167B parameters at 68\% of base) starts at ppl 1909. The conclusion is clear: conversion is not a free lunch at any tested tying level. Uptraining budget, not surgery, is the binding constraint, consistent with published recipes using 10--100B tokens. Practitioners should prioritize average initialization and reserve weight tying for the middle layers, where redundancy is highest.

The one-billion-parameter twin (Section~\ref{sec:results:twin}) fixes where the exchange rate stops being a safe recommendation, and the decision to halt there deserves its own statement because the closure contract would ordinarily have required two completed interventions first. Three facts decided it. Both arms score at chance on every public benchmark at 2.5B tokens, so the only signal the budget affords is held-out perplexity, and perplexity cannot tell a practitioner whether either model is worth serving. The tied arm's projected perplexity at 10B tokens, some seven days of the one available card, would leave it at chance still, while a Pythia-410m-class token budget lies roughly two hundred GPU-days away. And the interventions, whichever way they read, could only name the mechanism of a gap whose consequence for what ships was already settled. A programme whose objective was utility therefore stops at the boundary rather than spending its remaining compute refining a null; the record says so in those terms, declares no null, and leaves the mechanism as the first question for anyone who reopens the line.

The program closes with a sharply posed open problem: quality-preserving certification. The anytime-trained model (F24 B1) saturates quality at 0.991 parity but yields convergence rate 0.0 (solver uses all 12 iterations). The trajectory-penalized model (B2) achieves certified convergence (rate 1.0, 4.0 mean iterations) but at quality cost (0.577 BLiMP). No single arm meets both criteria, and their naive combination fails catastrophically (B4: 0.529). The mechanism is clear: under a fixed 12-iteration budget and the map's width-dependent contraction, quality and convergence pull against each other. Resolution requires either better fixed-point maps (higher contraction without capacity loss), adaptive iteration budgets that scale with width, or decoupled training for the two objectives---the solver-aware auxiliary loss (F16) shows one path, but composing it with anytime supervision remains open. This is the identified scientific frontier: not whether equilibrium language models can work, but how to certify them without sacrificing the quality gains that justify the approach.

